% !TeX root = ../../Book2.tex
\chapter{弱解与 Lax–Milgram 方法}
\label{b2:ch:15}
\BookChapterNumber{b2:ch:15}
% 本章摘要（不计入编号）
\begin{chapterabstract}
本章把散度型方程写成 Hilbert 空间中的测试等式。先明确 Sobolev 边界条件与负阶源项，再完整证明 Lax--Milgram 定理，并以 Poisson 方程说明存在性、唯一性、稳定性和有限维近似。一般有界可测系数无需具有经典导数；一致椭圆性控制主部，Gårding 不等式处理低阶项。最后借紧嵌入与第一卷的 Fredholm 理论，刻画原方程可能出现的有限维核及伴随相容条件。
\end{chapterabstract}

% 本章目标（不计入编号）
\begin{learninggoals}
\begin{itemize}
\item\label{b2:item:15-goal-weak}辨认测试空间所编码的内部方程与边界条件，正确使用 \(H^{-1}\) 数据。
\item\label{b2:item:15-goal-lax}分别核验型的有界性、强制性与空间完备性，掌握 Lax--Milgram 的闭值域证明。
\item\label{b2:item:15-goal-stability}用能量等式证明唯一性、比较和数据稳定性，并处理非齐次边界延拓。
\item\label{b2:item:15-goal-general}说明有界可测系数已足以给出弱存在性，区别原问题与移位问题。
\item\label{b2:item:15-goal-fredholm}在选读范围内，用伴随核写出 Fredholm 相容条件。
\end{itemize}
\end{learninggoals}

设一根杆的导热系数在接头处突然变化。平衡温度不应跳跃，穿过接头的热通量也应保持一致，但温度斜率可以不同。若只寻找处处二次可微的解，便会把这样的基本模型排除在外。弱形式保留可积梯度和通量平衡，把经典导数无法直接表达的问题放进一个完备的函数空间。

\ProseParagraphBreak

上一章用极小化构造对称能量的解；本章直接求解双线性等式，从而允许非对称系数。Hilbert 空间表示与紧算子的抽象结论采用第一卷第八章和附录 C 的证明，Sobolev 边界与紧嵌入采用第一卷第二十二至二十五章。相关理论的系统参考包括 Haim Brezis 的《Functional Analysis, Sobolev Spaces and Partial Differential Equations》\cite{Brezis2011Functional}与 Lawrence C. Evans 的《Partial Differential Equations》\cite{Evans2010Partial}。

% !TeX root = ../../Book2.tex
\section{弱形式}
\label{b2:sec:1501}

分段材料中的温度可以连续而斜率发生跳跃；此时逐点写二阶导数会掩盖界面上的通量条件。弱形式把微分移到测试函数上，使未知函数只需具有有限能量，同时仍保留方程及界面平衡。本节将第一章的局部弱解概念与上一章的 Sobolev 边界类结合起来，并为比 \(L^2\) 更粗糙的源项确定合适空间。

\subsection{测试函数}
\label{b2:subsec:150101}

在开集 \(\Omega\subset\mathbb R^d\) 中，测试函数取自 \(C_c^\infty(\Omega)\)。它们的支集离边界有正距离，因此分部积分不会引入尚未定义的边界值。若 \(u\in H_{\mathrm{loc}}^1(\Omega)\)、\(A\in L_{\mathrm{loc}}^\infty(\Omega;\mathbb R^{d\times d})\)，则
\[\langle-\operatorname{div}(A\nabla u),\varphi\rangle
 =\int_\Omega A\nabla u\cdot\nabla\varphi\,\mathrm dx\]
定义了一个分布。矩阵约定为 \((A\nabla u)_i=\sum_j a_{ij}\partial_j u\)，因而散度作用于通量的第一个指标。

\ProseParagraphBreak

局部测试函数足以辨认内部方程，但能量估计常须取 \(\varphi=u\)、截断后的 \(u\) 或差商。它们通常不光滑，也不紧支集。合法性依靠的是密度与连续延拓：若左右两边对 \(H_0^1\) 范数连续，而 \(C_c^\infty\) 在该空间稠密，则测试等式自动推广到整个 \(H_0^1\)。若只知道局部 \(H^1\)，便先乘紧支集截断函数，在较小区域完成同样步骤。不能在未建立这些连续性之前直接“取测试函数为解”。

\subsection{弱解}
\label{b2:subsec:150102}

以下主要讨论实函数，复函数可分别作用于实部与虚部。给定有界矩阵场 \(A\) 和源项 \(f\)，称 \(u\in H^1(\Omega)\) 满足散度方程的弱形式，是指
\[\int_\Omega A\nabla u\cdot\nabla\varphi=\langle f,\varphi\rangle
 \qquad(\varphi\in C_c^\infty(\Omega)).\]
这与第一章的定义一致；边界条件尚未包含在此等式中。如果 \(A,u\) 光滑到足以分部积分，它与经典方程等价；粗糙情形则以这个等式本身为准。

\begin{proposition}[一维界面的通量条件]
\label{b2:prop:15-interface-flux}
设 \(-1<s<1\)，\(a\) 在 \((-1,s)\)、\((s,1)\) 上分别为正常数，\(u\) 在两闭子区间上分别为 \(C^2\)，且在 \(s\) 处连续。若 \(-(au')'=f\) 在两侧经典成立，其中 \(f\in L^2(-1,1)\)，则 \(u\) 在整个区间为弱解，当且仅当
\[a(s-)u'(s-)=a(s+)u'(s+).\]
\end{proposition}
\begin{proof}
分段光滑且连续的 \(u\) 属于 \(H^1\)，其弱导数就是分段的一阶导数。在两段分别积分，外端因测试函数紧支集而消失，内端得到
\[\int_{-1}^1 au'\varphi'=\int_{-1}^1 f\varphi
  +\bigl(a(s-)u'(s-)-a(s+)u'(s+)\bigr)\varphi(s).\]
测试函数在 \(s\) 处可取任意值，因此等式对全部测试函数成立恰当且仅当括号为零。
\end{proof}

例如 \(a=1\) 于左半区间、\(a=2\) 于右半区间，令 \(u(x)=x\)（\(x\le0\)）、\(u(x)=x/2\)（\(x\ge0\)）。其导数跳跃，但通量恒为 \(1\)，所以 \(-(au')'=0\) 弱成立。若误要求 \(u'\) 连续，就会排除这个正确的平衡解；若允许 \(u\) 本身跳跃，又会产生分布导数中的点质量，通常已不属于所选能量空间。

\subsection{Sobolev 边界条件}
\label{b2:subsec:150103}

对任意有界开集，齐次 Dirichlet 条件定义为 \(u\in H_0^1(\Omega)\)；非齐次数据由一个延拓 \(g\in H^1(\Omega)\) 表示，条件为 \(u-g\in H_0^1(\Omega)\)。第一卷定理 23.1.6 说明，在有界 Lipschitz 区域上，这恰好等价于 Sobolev 迹满足 \(\operatorname{Tr}u=\operatorname{Tr}g\)。粗糙区域上仍可使用闭包定义，但不能未经证明把它等同于某个朴素逐点边界条件。

\ProseParagraphBreak

自然边界条件的写法不同。以 Poisson 方程为例，在 \(H^1(\Omega)\) 的全部方向测试时，光滑情形的 Green 公式是
\[\int_\Omega\nabla u\cdot\nabla v
  =\int_\Omega fv+\int_{\partial\Omega}g_N\operatorname{Tr}v\,\mathrm dS.\]
它编码 \(\partial_\nu u=g_N\)，而不是要求 \(u\) 在边界为零。取 \(v=1\) 立刻得到相容条件 \(\int_\Omega f+\int_{\partial\Omega}g_N=0\)。测试空间不同，编码的边界问题便不同；仅看内部方程无法辨认二者。

\ProseParagraphBreak

即使区域光滑，也不宜把任意边界函数都当作合法 \(H^1\) Dirichlet 数据。数据应当处于迹空间，并存在 \(H^1\) 延拓；第一卷第二十三章在 Lipschitz 图中用分数阶范数刻画这种要求。本章直接给定延拓 \(g\)，从而把数据的合法性写进假设，不需要在每个存在性证明里重新构造延拓算子。

\subsection{负阶 Sobolev 空间初步}
\label{b2:subsec:150104}

设 \(\Omega\) 有界，以梯度范数赋予 \(H_0^1(\Omega)\) Hilbert 结构。定义
\[H^{-1}(\Omega)=(H_0^1(\Omega))^*,\qquad
 \lVert f\rVert_{H^{-1}}=\sup_{\lVert\nabla v\rVert_2\le1}|\langle f,v\rangle|.
\]
这是本章使用的 \term[fujie-sobolev-kongjian]{负阶 Sobolev 空间}[negative-order Sobolev space]的第一种情形。换成全 \(H^1\) 范数会得到等价范数；双空间本身不变。由 Poincaré 不等式，\(L^2\) 连续嵌入 \(H^{-1}\)。更一般地，若 \(F\in L^2(\Omega;\mathbb R^d)\)，则 \(-\operatorname{div}F\in H^{-1}\)，且其范数不超过 \(\lVert F\rVert_2\)。

\begin{proposition}[负阶数据的通量表示]
\label{b2:prop:15-negative-representation}
对每个 \(f\in H^{-1}(\Omega)\)，存在 \(F\in L^2(\Omega;\mathbb R^d)\)，使
\[\langle f,v\rangle=\int_\Omega F\cdot\nabla v\,\mathrm dx
 \quad(v\in H_0^1(\Omega)),\qquad \lVert F\rVert_2=\lVert f\rVert_{H^{-1}}.\]
\end{proposition}
\begin{proof}
映射 \(v\mapsto\nabla v\) 把 \(H_0^1\) 等距嵌入 \(L^2(\Omega;\mathbb R^d)\)，其像因完备性而闭。在线性子空间上定义 \(\Lambda(\nabla v)=\langle f,v\rangle\)；Poincaré 不等式保证定义无歧义。Hahn--Banach 定理将其等范数延拓到整个向量值 \(L^2\)。第一卷定理 8.4.7 的 Hilbert 空间 Riesz 表示给出向量 \(F\)，其范数等于泛函范数，恰得所列公式。
\end{proof}

这种表示未必唯一，因为与所有梯度正交的通量可以自由加到 \(F\) 上。负阶空间描述的是通量的散度，而不是某个唯一的通量。它也确实容纳比函数更粗糙的数据：在一维有界区间内，点值由 \(H_0^1\) 范数控制，故点质量属于 \(H^{-1}\)；在二维及更高维，一般不再如此。后面的习题用缩放说明这一维数差异，避免把一维直觉直接移到高维。

% !TeX root = ../../Book2.tex
\section{Hilbert 空间方法}
\label{b2:sec:1502}

上一章的极小化要求能量保留方程的全部信息。对非对称系数，这个要求可能失败；双线性型方法则直接研究等式 \(a(u,v)=\ell(v)\)。其中未知量占第一个位置，测试向量占第二个位置。本节把有界性和强制性分别转成算子的上界和下界，再用正交性证明可解性。

\subsection{双线性型}
\label{b2:subsec:150201}

设 \(H\) 是实 Hilbert 空间。\term[shuangxianxing-xing]{双线性型}[bilinear form]是对两个变量分别线性的映射 \(a:H\times H\to\mathbb R\)。对称性 \(a(u,v)=a(v,u)\) 是额外要求，本章的主要存在定理不作这一假设。例如
\[a(u,v)=\int_\Omega A\nabla u\cdot\nabla v\,\mathrm dx\]
在 \(H_0^1(\Omega)\) 上是双线性型；若矩阵 \(A(x)\) 对称，则型对称。一个纯反对称矩阵场对 \(a(u,u)\) 没有贡献，却可能改变 \(a(u,v)\)。

\ProseParagraphBreak

复 Hilbert 空间中通常改用对第一个变量线性、对第二个变量共轭线性的半双线性型。相应右端也对测试向量共轭线性，强制性写为 \(\operatorname{Re}a(u,u)\ge\alpha\lVert u\rVert^2\)。这与本书内积对第一变量线性的约定一致。为使证明中的对偶方向透明，下面先给出实空间版本，随后说明复版本的调整；不要把这一约定与分布配对的双线性约定混用。

\subsection{有界性}
\label{b2:subsec:150202}

双线性型的有界性是存在 \(M<\infty\)，使
\[|a(u,v)|\le M\lVert u\rVert_H\lVert v\rVert_H\qquad(u,v\in H).\]
它保证固定任一变量后所得泛函连续，也保证从稠密的测试子空间推广等式时极限合法。若 \(A\) 的矩阵算子范数几乎处处不超过 \(\Lambda\)，则
\[|a(u,v)|\le\Lambda\lVert\nabla u\rVert_2\lVert\nabla v\rVert_2.\]
这里有界性只使用系数的本性上界，不涉及其导数，也不要求点点连续。

\begin{proposition}[有界型的算子表示]
\label{b2:prop:15-form-operator}
设 \(a\) 是实 Hilbert 空间 \(H\) 上的有界双线性型。则存在唯一有界线性算子 \(A:H\to H\)，使
\[a(u,v)=\langle Au,v\rangle_H\qquad(u,v\in H),\]
且 \(\lVert A\rVert=\sup_{\lVert u\rVert,\lVert v\rVert\le1}|a(u,v)|\)。型对称当且仅当算子自伴。
\end{proposition}
\begin{proof}
固定 \(u\)，第一卷定理 8.4.7 把连续线性泛函 \(v\mapsto a(u,v)\) 表示为与唯一向量 \(Au\) 的内积。表示的唯一性和第一变量线性性给出 \(A\) 的线性性，且 \(\lVert Au\rVert\le M\lVert u\rVert\)。对单位向量分别取上确界给出范数等式。最后，\(a(u,v)=a(v,u)\) 等价于 \(\langle Au,v\rangle=\langle u,Av\rangle\)，正是自伴条件。
\end{proof}

这个 \(A\) 是依赖所选 Hilbert 内积的表示算子，不必等于微分表达式本身。例如在梯度内积的 \(H_0^1\) 中，Poisson 型 \(\int\nabla u\cdot\nabla v\) 的表示算子就是恒等算子；作为分布的 \(-\Delta\) 则把 \(H_0^1\) 映到 \(H^{-1}\)。分清值域有助于避免把有界算子与无界微分算子的谱性质混在一起。

\subsection{强制性}
\label{b2:subsec:150203}

称型 \(a\) 强制，是指存在 \(\alpha>0\)，使
\[a(u,u)\ge\alpha\lVert u\rVert_H^2\qquad(u\in H).\]
这里要求统一的 \(\alpha\)，而不只是每个非零向量上取正值。结合算子表示与 Cauchy--Schwarz 不等式，立刻得到
\[\alpha\lVert u\rVert_H\le\lVert Au\rVert_H.\]
它同时给出单射性和逆映射在值域上的有界性。不过，单有这种下界并不保证一般算子满射；例如 \(\ell^2\) 上的右移是等距单射却漏掉第一个坐标方向。Lax--Milgram 证明还要利用同一个型在测试向量处的正性，排除值域的正交补。

\ProseParagraphBreak

对椭圆型矩阵，条件 \(\xi^{\mathsf T}A(x)\xi\ge\lambda|\xi|^2\) 几乎处处成立时，梯度型在 \(H_0^1\) 的梯度范数下以 \(\alpha=\lambda\) 强制。在整个 \(H^1\) 上却并非如此，因为常函数在型中为零。选空间与选范数是定理假设的一部分，不能在写下“椭圆”之后省略核验。

\subsection{Lax–Milgram 定理}
\label{b2:subsec:150204}

\begin{theorem}[Lax--Milgram]
\label{b2:thm:15-lax-milgram}
设 \(H\) 是实 Hilbert 空间，\(a:H\times H\to\mathbb R\) 是双线性型，且对某些 \(M<\infty\)、\(\alpha>0\) 有
\[|a(u,v)|\le M\lVert u\rVert_H\lVert v\rVert_H,
 \qquad a(u,u)\ge\alpha\lVert u\rVert_H^2.\]
则对每个 \(\ell\in H^*\)，存在唯一 \(u\in H\) 满足
\[a(u,v)=\ell(v)\qquad(v\in H),\]
且 \(\lVert u\rVert_H\le\alpha^{-1}\lVert\ell\rVert_{H^*}\)。
\end{theorem}
\begin{proof}
用 \cref{b2:prop:15-form-operator} 表示 \(a(u,v)=\langle Au,v\rangle\)，并用 Riesz 表示 \(\ell(v)=\langle f,v\rangle\)。所求等式等价于 \(Au=f\)。强制性给出 \(\lVert Au\rVert\ge\alpha\lVert u\rVert\)，所以 \(A\) 单射。

若 \(Au_n\to y\)，下界说明 \((u_n)\) 是 Cauchy 序列，因而收敛到某个 \(u\in H\)。由 \(A\) 连续，\(Au=y\)，故值域闭。若 \(w\) 与值域正交，则 \(a(v,w)=\langle Av,w\rangle=0\) 对所有 \(v\in H\) 成立。取 \(v=w\)，强制性推出 \(w=0\)。因此值域的正交补为零；闭子空间的正交分解给出值域等于 \(H\)。于是 \(Au=f\) 唯一可解，且算子下界给出 \(\lVert u\rVert\le\alpha^{-1}\lVert f\rVert=\alpha^{-1}\lVert\ell\rVert\)。
\end{proof}

定理也称 \term[lax-milgram-dingli]{Lax--Milgram 定理}[Lax--Milgram theorem]。复空间版本把强制性改成实部下界，并把 \(\ell\) 视为连续共轭线性泛函；算子表示、闭值域和正交补的三步证明保持不变，在取 \(v=w\) 时使用 \(\operatorname{Re}a(w,w)\) 即可。

\ProseParagraphBreak

还可以把存在性看成一次收缩迭代。若 \(A\) 满足上述上下界，取 \(0<t<2\alpha/M^2\)，则
\[
 \lVert(I-tA)w\rVert^2
 =\lVert w\rVert^2-2t\langle Aw,w\rangle+t^2\lVert Aw\rVert^2
 \le(1-2t\alpha+t^2M^2)\lVert w\rVert^2.
\]
因此 \(u_{n+1}=u_n+t(f-Au_n)\) 收敛到唯一解。这个计算不仅提供另一种构造，也表明当 \(\alpha\) 很小而 \(M\) 很大时，允许的步长和收敛速度都会变差。定理的稳定性常数真实反映了问题接近失去强制性的程度。

% !TeX root = ../../Book2.tex
\section{Poisson 方程}
\label{b2:sec:1503}

现在把抽象定理用于最基本的边值问题。区域只需有界，源项只需属于 \(H^{-1}\)；无需显式 Green 函数，也无需预先知道解的二阶导数。存在性、稳定性与边界数据的处理都在同一个能量空间中完成。

\subsection{弱解存在性}
\label{b2:subsec:150301}

\begin{theorem}[Dirichlet 问题的弱可解性]
\label{b2:thm:15-poisson-existence}
设 \(\Omega\subset\mathbb R^d\) 为有界开集，\(f\in H^{-1}(\Omega)\)。则存在唯一 \(u\in H_0^1(\Omega)\)，使
\[\int_\Omega\nabla u\cdot\nabla v\,\mathrm dx=\langle f,v\rangle
 \qquad(v\in H_0^1(\Omega)).\]
并且 \(\lVert\nabla u\rVert_2\le\lVert f\rVert_{H^{-1}}\)。
\end{theorem}
\begin{proof}
取 \(H=H_0^1(\Omega)\)，内积为 \(\langle u,v\rangle_H=\int\nabla u\cdot\nabla v\)。由 Poincaré 不等式，它与完整 Sobolev 范数等价，故 \(H\) 完备。型 \(a(u,v)=\langle u,v\rangle_H\) 的有界常数与强制常数均为 \(1\)，右端按 \(H^{-1}\) 定义连续。\cref{b2:thm:15-lax-milgram} 因而给出存在、唯一性及估计。
\end{proof}

本例中 Lax--Milgram 退化为 Riesz 表示；一般系数时，同一个定理才显出更大的作用。该弱解也恰为 \(\tfrac12\lVert\nabla v\rVert_2^2-\langle f,v\rangle\) 的唯一极小元：只要将上一章的二次展开中的 \(L^2\) 配对换成对偶配对，证明逐字成立。两种构造在对称情形相遇，但弱型方法还可以继续处理非对称问题。

\ProseParagraphBreak

对 \(\Omega=(0,1)\)、\(f=\delta_s\)、\(0<s<1\)，解为
\[
 u(x)=\begin{cases}(1-s)x,&0\le x\le s,\\s(1-x),&s\le x\le1.\end{cases}
\]
它连续、分段仿射，导数在 \(s\) 处向下跳 \(1\)。两段分部积分给出 \(\int_0^1u'v'=v(s)\)，先对光滑 \(v\) 成立，再由点值的连续性推广。因此点源产生一个合法弱解，却不产生经典 \(C^2\) 解；这不是存在性定理的缺陷，而是数据本身的粗糙性。

\subsection{唯一性}
\label{b2:subsec:150302}

唯一性的能量论证值得单独保留，因为它不依赖解是如何构造的。若 \(u_1,u_2\) 具有同一数据，差 \(w\in H_0^1\) 满足 \(\int\nabla w\cdot\nabla v=0\)。取 \(v=w\)，得到 \(\nabla w=0\)；Poincaré 不等式推出 \(w=0\)。这个结论对不连通区域也成立，因为零边界约束在每个分支上排除了常数自由度。

\ProseParagraphBreak

内部方程则没有唯一性。例如任一调和函数都可加到 Poisson 方程的一个局部解上，而不改变源项。纯 Neumann 问题即使规定通量，仍允许每个连通分支上的加法常数。因而“取差再测试”最后得出梯度为零之后，还必须指出究竟由哪项边界条件或归一化排除了常数，不能直接省略最后一步。

\begin{proposition}[弱比较原理]
\label{b2:prop:15-weak-comparison}
设 \(u_1,u_2\in H_0^1(\Omega)\) 分别解 \(-\Delta u_i=f_i\)，且 \(\langle f_1-f_2,v\rangle\le0\) 对每个非负 \(v\in H_0^1(\Omega)\) 成立。则 \(u_1\le u_2\) 几乎处处。
\end{proposition}
\begin{proof}
令 \(w=u_1-u_2\)。Sobolev 截断性质保证 \(w_+=\max\{w,0\}\in H_0^1\)，且 \(\nabla w_+=\mathbf1_{\{w>0\}}\nabla w\)。以它测试差方程，得
\[\lVert\nabla w_+\rVert_2^2=\langle f_1-f_2,w_+\rangle\le0.\]
故 \(w_+=0\)，即所需结论。截断性质可先用光滑 Lipschitz 逼近正部函数、链式法则和 \(H^1\) 极限得到，再对 \(H_0^1\) 使用零边界逼近；这里没有要求解逐点可微。
\end{proof}

特别地，非负源项与零边界数据给出非负弱解，与第十三章经典最大值原理的符号一致。比较的结论是几乎处处关系；在尚未得到连续代表元之前，不应改说“每个点都非负”。

\subsection{能量估计}
\label{b2:subsec:150303}

取解自身作测试，Cauchy--Schwarz 的对偶形式给出
\[\lVert\nabla u\rVert_2^2=\langle f,u\rangle
 \le\lVert f\rVert_{H^{-1}}\lVert\nabla u\rVert_2.\]
若梯度非零可约去一因子；若为零估计显然成立。因此解映射 \(H^{-1}\to H_0^1\) 连续。对于两组数据，同一计算给出
\[\lVert\nabla(u_1-u_2)\rVert_2\le\lVert f_1-f_2\rVert_{H^{-1}}.\]
特别地，平滑数据的近似会在能量范数中产生相应解的收敛，这为从光滑近似过渡到粗糙数据提供了稳定性依据。

\ProseParagraphBreak

若 \(f\in L^2\)，进一步有 \(\lVert\nabla u\rVert_2\le C_P\lVert f\rVert_2\)、\(\lVert u\rVert_2\le C_P^2\lVert f\rVert_2\)。注意区域放大时 \(C_P\) 也会变化：在 \(r\Omega\) 上换元表明 \(C_P\) 按长度尺度增长。这与二阶方程解相对于源项多出两个长度因子的事实一致，估计中的几何依赖不能无故删除。

\begin{proposition}[有限维近似的最佳误差]
\label{b2:prop:15-galerkin-poisson}
设 \(V\subset H_0^1(\Omega)\) 是有限维子空间，\(u\) 为 \cref{b2:thm:15-poisson-existence} 的解。则存在唯一 \(u_V\in V\) 满足 \(\int\nabla u_V\cdot\nabla v=\langle f,v\rangle\) 对所有 \(v\in V\) 成立，而且
\[\lVert\nabla(u-u_V)\rVert_2=\inf_{v\in V}\lVert\nabla(u-v)\rVert_2.\]
\end{proposition}
\begin{proof}
限制到 \(V\) 后型仍有界、强制，故有限维方程唯一可解。两式相减得 \(\int\nabla(u-u_V)\cdot\nabla v=0\) 对所有 \(v\in V\) 成立。对任意 \(v\in V\)，分解 \(u-v=(u-u_V)+(u_V-v)\)，两项在梯度内积下正交，故误差平方等于两项范数平方之和。最小值恰在 \(v=u_V\) 达到。
\end{proof}

这个结论不宣称任一有限维空间都给出好的逼近；它说解离该空间有多近，近似解就有多好。若一列子空间递增且其并在 \(H_0^1\) 中稠密，最佳误差便趋零。递增性可改成对每个固定函数都能沿整个序列逼近的条件，但只有并集稠密不够，例如在良好空间之间反复插入零空间便会破坏收敛。存在性、离散方程的可解性与近似空间的逼近能力仍是三个可以分别核验的事实。

\ProseParagraphBreak

在区间 \((0,1)\) 上，这个结论可以算得很具体。取均匀节点 \(x_j=jh\)、\(h=1/N\)，令 \(V_h\) 为节点间分段仿射、两端为零的连续函数空间。基函数 \(\phi_j\) 在 \(x_j\) 取值一，在其余节点取值零。每个基函数只有相邻两段上的斜率非零，直接积分得
\[\int_0^1\phi_j'^2=\frac2h,\qquad
 \int_0^1\phi_j'\phi_{j+1}'=-\frac1h,\qquad
 \int_0^1\phi_j'\phi_k'=0\quad(|j-k|>1).\]
因此弱解的有限维近似由三对角线性方程组确定。矩阵的正定性并不需要另作行列式计算：任意系数向量对应的二次型就是分段仿射函数的梯度平方积分，零值只可能对应零函数。

\ProseParagraphBreak

例如取 \(f=1\)，精确弱解为 \(u(x)=x(1-x)/2\)。它在各小区间的导数平均等于两端函数值之差除以 \(h\)，而任意 \(v_h\in V_h\) 的导数在该段为常数。故 \(u\) 的节点线性插值 \(I_hu\) 满足 \(\int(u'- (I_hu)')v_h'=0\)，由唯一性即知 \(u_h=I_hu\)。在长度 \(h\) 的一段上，导数误差是以中点为零的线性函数，其平方积分为 \(h^3/12\)；函数误差则为 \((x-x_j)(x_{j+1}-x)/2\)。逐段积分并相加得到
\[\lVert u'-u_h'\rVert_2=\frac h{\sqrt{12}},\qquad
 \lVert u-u_h\rVert_2=\frac{h^2}{\sqrt{120}}.\]
这里第二个收敛速度比第一个快，来自二次函数的具体结构，不能仅从抽象存在定理读出。一般数据下的误差阶还需要解的正则性和近似空间的逼近估计；下一章提供其中的正则性工具。

\ProseParagraphBreak

弱解的误差也可以用残差直接测量。对任意 \(w\in H_0^1\) 定义 \(r=f+\Delta w\in H^{-1}\)，则 \(\langle r,v\rangle=\int\nabla(u-w)\cdot\nabla v\)。对单位梯度范数的测试函数取上确界，Cauchy--Schwarz 给出上界，而在误差非零时取 \(v=(u-w)/\lVert\nabla(u-w)\rVert_2\) 达到该界。因此 \(\lVert r\rVert_{H^{-1}}=\lVert\nabla(u-w)\rVert_2\)。对分段仿射近似，残差可能含节点通量跳跃，不能只计算每段内为零的二阶导数；负范数恰好保留这些分布贡献。

\subsection{非齐次边界数据}
\label{b2:subsec:150304}

给定 \(g\in H^1(\Omega)\)，令 \(u=g+w\)，其中 \(w\in H_0^1\)。方程等价于
\[\int_\Omega\nabla w\cdot\nabla v
  =\langle f,v\rangle-\int_\Omega\nabla g\cdot\nabla v\qquad(v\in H_0^1).\]
右端属于 \(H^{-1}\)，范数不超过 \(\lVert f\rVert_{H^{-1}}+\lVert\nabla g\rVert_2\)。因此齐次理论给出唯一 \(w\)，从而给出唯一满足 \(u-g\in H_0^1\) 的解；并有
\[\lVert\nabla w\rVert_2\le\lVert f\rVert_{H^{-1}}+\lVert\nabla g\rVert_2.\]
完整 \(H^1\) 估计还要加上 \(g\) 的 \(L^2\) 范数。

\ProseParagraphBreak

若另一个延拓 \(\widetilde g\) 与 \(g\) 相差 \(H_0^1\) 函数，两者给出的仿射可行集完全相同；由唯一性得到同一个 \(u\)。所以解依赖边界类，而不依赖人为选择的延拓。若要比较两组边界数据，先选择它们差的一个延拓，再把上面的稳定性估计应用于差方程；优化延拓的范数便自然引出商空间形式的迹范数。

\ProseParagraphBreak

例如在单位球上给定仿射边界数据 \(g(x)=a\cdot x+b\)，且 \(f=0\)。由于 \(\int a\cdot\nabla v=0\) 对 \(v\in H_0^1\) 成立，齐次修正为零，故 \(u=g\)。这里无需使用球的 Poisson 核；而对一般边界数据，弱方法仍给出存在性，之后才可能借第十三章的表示公式识别它。不同构造的解一旦处于同一能量边界类，就由唯一性自动一致。

% !TeX root = ../../Book2.tex
\section{一般二阶椭圆算子}
\label{b2:sec:1504}

Poisson 方程的证明只用到梯度型的有界性与正性，因而可以推广到系数随位置变化的算子。困难主要来自两处：低阶项可能破坏正性，非对称主部使极小化不再适用。本节先得到移位后的可解性，再在选读小节中借紧性解释原问题的有限维障碍。

\subsection{散度型算子}
\label{b2:subsec:150401}

设 \(\Omega\subset\mathbb R^d\) 有界，考虑
\[Lu=-\operatorname{div}(A\nabla u)+b\cdot\nabla u+cu,\]
其中矩阵 \(A\)、向量 \(b\) 和实函数 \(c\) 均有界可测。其 \term[sandu-xing-suanzi]{散度型算子}[divergence-form operator]主部用通量 \(A\nabla u\) 定义，不必对系数求经典导数。对应型为
\[a(u,v)=\int_\Omega A\nabla u\cdot\nabla v
       +\int_\Omega(b\cdot\nabla u)v+\int_\Omega cuv.
\]
弱 Dirichlet 问题是在 \(H_0^1(\Omega)\) 中求 \(u\)，使 \(a(u,v)=\langle f,v\rangle\) 对所有 \(v\in H_0^1\) 成立。

\ProseParagraphBreak

假设存在 \(\lambda>0\)，使对几乎处处的 \(x\) 和全部实向量 \(\xi\) 有
\[\xi^{\mathsf T}A(x)\xi\ge\lambda|\xi|^2.\]
这称为一致椭圆性。由于反对称部分的二次型为零，它等价于对称部分的相同下界；但方程本身仍保留全部矩阵。若 \(A\) 不可微，不能把 \(-\partial_i(a_{ij}\partial_j u)\) 改成 \(-a_{ij}\partial_{ij}u\) 并丢掉系数导数，这两种形式是不同的方程类别。

\subsection{有界可测系数}
\label{b2:subsec:150402}

设 \(\lVert A(x)\rVert\le\Lambda\)、\(\lVert b\rVert_\infty=B\)、\(\lVert c\rVert_\infty=C_0\)。逐项应用 Hölder 与 Poincaré 不等式，得到
\[|a(u,v)|\le(\Lambda+BC_P+C_0C_P^2)\lVert\nabla u\rVert_2\lVert\nabla v\rVert_2.\]
所以有界可测系数已经足够定义连续弱型。若 \(b=0\)、\(c\ge0\)，椭圆性还给出 \(a(u,u)\ge\lambda\lVert\nabla u\rVert_2^2\)。由 Lax--Milgram，对每个 \(f\in H^{-1}\) 唯一可解，且
\[\lVert\nabla u\rVert_2\le\lambda^{-1}\lVert f\rVert_{H^{-1}}.\]
这些结论没有要求矩阵对称。

\ProseParagraphBreak

系数粗糙时，解的导数可能在界面跳跃，\cref{b2:prop:15-interface-flux} 已给出最简单的例子。这说明存在性定理与下一章的 \(H^2\) 正则性定理必须使用不同假设：前者只要系数有界可测，后者的差商证明需要系数具有 Lipschitz 控制。不能因为方程是“椭圆的”，就把所有正则性结论一并附加上去。

\ProseParagraphBreak

对于非齐次边界数据，仍写 \(u=g+w\)。右端变成 \(f-a(g,\cdot)\)；由同样的积分估计，它是 \(H_0^1\) 上的连续泛函。因此齐次问题的结论立即推广。需要注意，估计 \(a(g,v)\) 时 \(g\) 不属于 \(H_0^1\)，不能对 \(g\) 使用零边界 Poincaré 不等式，而应保留其完整 \(H^1\) 范数。

\subsection{Gårding 不等式}
\label{b2:subsec:150403}

\begin{proposition}[Gårding 不等式]
\label{b2:prop:15-garding}
设 \(\Omega\) 有界，\(A,b,c\) 如上有界可测，且 \(A\) 以常数 \(\lambda>0\) 一致椭圆。令
\[\beta=\frac{\lVert b\rVert_\infty^2}{2\lambda}+\lVert c_-\rVert_\infty,
 \qquad c_- =\max\{-c,0\}.\]
则对所有 \(u\in H_0^1(\Omega)\)，
\[a(u,u)\ge\frac\lambda2\lVert\nabla u\rVert_2^2-\beta\lVert u\rVert_2^2.\]
因而 \(a_\mu(u,v)=a(u,v)+\mu\int_\Omega uv\) 在每个 \(\mu\ge\beta\) 时强制；方程 \((L+\mu I)u=f\) 对每个 \(f\in H^{-1}\) 唯一弱可解。
\end{proposition}
\begin{proof}
主部下界是 \(\lambda\lVert\nabla u\rVert_2^2\)。一阶项由 Young 不等式满足
\[\left|\int(b\cdot\nabla u)u\right|
 \le B\lVert\nabla u\rVert_2\lVert u\rVert_2
 \le\frac\lambda2\lVert\nabla u\rVert_2^2+\frac{B^2}{2\lambda}\lVert u\rVert_2^2.\]
零阶项下界为 \(-\lVert c_-\rVert_\infty\lVert u\rVert_2^2\)。相加即得不等式；加上 \(\mu\lVert u\rVert_2^2\) 后，强制常数至少为 \(\lambda/2\)。有界性已证明，故应用 \cref{b2:thm:15-lax-milgram} 得到最后结论。
\end{proof}

上述估计称为 \term[garding-budengshi]{Gårding 不等式}[Gårding inequality]。它不是原型的强制性，而是允许减去一个较低阶的 \(L^2\) 项。移位恢复正性，却也改变了方程：解 \((L+\mu I)u=f\) 并不能直接宣称已经解了 \(Lu=f\)。去掉移位需要额外论证，下一小节用紧嵌入完成这一步。

\ProseParagraphBreak

有时漂移项可以更精确处理。若 \(b\) 足够光滑，且 \(u\in H_0^1\)，则密度和分部积分给出 \(\int(b\cdot\nabla u)u=-\tfrac12\int(\operatorname{div}b)u^2\)。若 \(c-\tfrac12\operatorname{div}b\ge0\)，原型已强制，无须移位。对于仅有界可测的 \(b\)，散度未必是可比较的函数，此时前面的 Young 估计才是无需额外正则性的可靠结论。

\OptionalSubsection{Fredholm 择一定理}{b2:subsec:150404}

本小节只使用第一卷已经完整证明的紧算子 Fredholm 理论。应用的紧性来自第一卷推论 25.2.5：任意有界开集上的 \(H_0^1(\Omega)\hookrightarrow L^2(\Omega)\) 紧嵌入。它不需要 Lipschitz 边界，因为零延拓与平移紧性已经控制了边界附近的质量。

\begin{theorem}[椭圆边值问题的有限维障碍]
\label{b2:thm:15-fredholm-elliptic}
在 \cref{b2:prop:15-garding} 的假设下，令
\[
 \begin{aligned}
 N&=\{u\in H_0^1:a(u,v)=0\ \text{对所有 }v\in H_0^1\},\\
 N^*&=\{w\in H_0^1:a(v,w)=0\ \text{对所有 }v\in H_0^1\}.
 \end{aligned}
\]
则两空间有限维且维数相等。对 \(f\in H^{-1}\)，弱问题 \(a(u,v)=\langle f,v\rangle\) 有解，当且仅当
\[\langle f,w\rangle=0\qquad(w\in N^*).\]
有解时全部解为 \(u_0+N\)。特别地，齐次问题只有零解，当且仅当任意右端都有唯一解。
\end{theorem}
\begin{proof}
记 \(H=H_0^1\)，以梯度为内积。分别用算子 \(T,T_\mu,M:H\to H\) 表示型 \(a,a_\mu\) 及 \(\int uv\)。选 \(\mu\ge\beta\)，Lax--Milgram 使 \(T_\mu\) 成为有界可逆算子。又 \(T=T_\mu-\mu M=T_\mu(I-K)\)，其中 \(K=\mu T_\mu^{-1}M\)。

为证 \(M\) 紧，取 \(H\) 中有界列 \((u_n)\)。紧嵌入给出在 \(L^2\) 中收敛的子列；且
\[\lVert M(u_n-u_k)\rVert_H
 =\sup_{\lVert v\rVert_H\le1}\left|\int(u_n-u_k)v\right|
 \le C_P\lVert u_n-u_k\rVert_2,\]
所以像子列在 \(H\) 中收敛。由此 \(K\) 紧。第一卷定理 C.4.4 说明 \(I-K\) 有闭值域、有限维核与余核，二者维数相等；左乘可逆算子后，\(T\) 保留这些性质。

在 Hilbert 空间中，闭值域满足 \(\operatorname{ran}T=(\ker T^*)^\perp\)。而 \(\ker T=N\)、\(\ker T^*=N^*\)，因为 \(\langle Tv,w\rangle_H=a(v,w)\)。将右端 \(f\) 用 Riesz 向量表示，可解性即等价于它与 \(N^*\) 正交，也就是所列对偶配对为零。两个解之差恰在 \(N\) 内，得到所有结论。
\end{proof}

对称型的 \(N^*\) 就是 \(N\)；非对称型必须检查伴随核，不能用原核替代。例如有限维非对称矩阵的值域正交于转置矩阵的核，而未必正交于自身的核，这里的区别完全相同。定理中的相容条件只有有限多个，可以对 \(N^*\) 的一组基逐个检验。

\ProseParagraphBreak

在 \((0,\pi)\) 上取 \(L=-\frac{\mathrm d^2}{\mathrm dx^2}-1\) 和零端点条件。齐次核为 \(\operatorname{span}\{\sin x\}\)，故 \(Lu=f\) 可解恰当且仅当 \(\int_0^\pi f(x)\sin x\,\mathrm dx=0\)。源项 \(f=\sin x\) 因此不可解，而 \(f=\sin2x\) 有解 \(u=\frac13\sin2x+c\sin x\)。这个例子同时显示：一致椭圆主部并不排除共振；移位后的唯一可解性与原问题的有限维障碍可以同时成立。

% 本章小结（不计入编号）
\begin{chaptersummary}
弱形式减少了未知函数所需的导数阶数，却没有丢掉方程：紧支集测试辨认内部散度，测试空间和边界类决定所解的边值问题，界面项记录通量平衡。负阶空间把函数源项与散度型源项放在同一连续对偶中。Lax--Milgram 用有界性构造算子，以强制性控制逆，再排除值域的正交补，得到无需对称性的存在与稳定性。

\ProseParagraphBreak

一般椭圆主部只需有界可测系数即可定义能量弱解；低阶项可能使原型失去强制性，但加上足够大的零阶移位可恢复它。去掉移位时，紧嵌入将问题化为恒等算子的紧扰动，可解性由伴随核的有限个条件决定。整个论证首先产生 \(H^1\) 解；要得到 \(H^2\)、连续梯度乃至经典解，下一章还须逐项增加系数、数据和边界的正则性。
\end{chaptersummary}


\BookExercises

各题均为本书自拟。前几题核验空间和数据，随后比较近似、系数扰动与共振；最后一题使用选读小节的相容条件。

% 自拟：一维点源的精确对偶范数。
\begin{exercise}\label{b2:ex:15-point-source}
在 \((0,1)\) 上，以梯度范数定义 \(H^{-1}\)。对 \(0<s<1\)，求点质量 \(\delta_s\) 的精确范数，并说明哪个函数达到对偶上确界。
\end{exercise}
\begin{solution}
令 \(G_s(x)=(1-s)x\)（\(x\le s\)）、\(G_s(x)=s(1-x)\)（\(x\ge s\)）。正文已验证 \(v(s)=\int G_s'v'\)。因此 \(|v(s)|\le\lVert G_s'\rVert_2\lVert v'\rVert_2\)，且取 \(v=G_s/\lVert G_s'\rVert_2\) 达到上界。直接计算 \(\lVert G_s'\rVert_2^2=s(1-s)^2+(1-s)s^2=s(1-s)\)，所以范数为 \(\sqrt{s(1-s)}\)。
\end{solution}

% 自拟：高维点质量的尺度障碍。
\begin{exercise}\label{b2:ex:15-high-dimensional-point}
设 \(B(0,1)\subset\Omega\subset\mathbb R^d\) 有界。证明当 \(d\ge2\) 时点质量 \(\delta_0\notin H^{-1}(\Omega)\)。二维可使用在 \(r\le\varepsilon\) 等于 \(1\)、在 \(\varepsilon<r<1\) 等于 \(\log r/\log\varepsilon\)、在 \(r\ge1\) 等于零的径向截断。
\end{exercise}
\begin{solution}
若点质量属于该对偶，应有 \(|\varphi(0)|\le C\lVert\nabla\varphi\rVert_2\) 对所有光滑紧支集函数成立。\(d>2\) 时取 \(\varphi_\varepsilon(x)=\varphi(x/\varepsilon)\)，其中 \(\varphi(0)=1\)、支集在单位球；梯度范数平方为 \(\varepsilon^{d-2}\lVert\nabla\varphi\rVert_2^2\to0\)，矛盾。二维的题设截断 \(w_\varepsilon\) 有 \(w_\varepsilon(0)=1\)，且
\[\int|\nabla w_\varepsilon|^2=\frac{2\pi}{|\log\varepsilon|^2}\int_\varepsilon^1\frac{\mathrm dr}{r}=\frac{2\pi}{|\log\varepsilon|}\to0.\]
可先把外半径换成固定小于 \(1\) 的数，再用半径小于 \(\varepsilon/2\) 的磨光核平滑。中心仍等于 \(1\)，支集仍紧含于单位球，梯度范数不增；故得到同样矛盾。二维必须使用对数截断，简单缩放的梯度范数恰好不变。
\end{solution}

% 自拟：分段材料的显式有效导热。
\begin{exercise}\label{b2:ex:15-layered-material}
令 \(a=1\) 于 \((0,1/2)\)、\(a=4\) 于 \((1/2,1)\)。求 \(-(au')'=0\)、\(u(0)=0,u(1)=1\) 的弱解，计算左右斜率及总能量，并判断它是否属于 \(H^2(0,1)\)。
\end{exercise}
\begin{solution}
通量为常数 \(q\)，且 \(1=\int_0^1u'=q(1/2+1/8)=5q/8\)，所以 \(q=8/5\)。左斜率为 \(8/5\)，右斜率为 \(2/5\)，解是通过 \((0,0),(1/2,4/5),(1,1)\) 的分段仿射函数。能量 \(\tfrac12\int a|u'|^2=\tfrac12q\int u'=4/5\)。一阶导数在中点有非零跳跃，二阶分布导数含点质量，故不属于 \(L^2\)，解不在 \(H^2\)。
\end{solution}

% 自拟：一般强制型的Galerkin误差而非对称投影。
\begin{exercise}\label{b2:ex:15-cea}
设 \(a\) 满足 Lax--Milgram 条件，常数为 \(M,\alpha\)，\(V\subset H\) 有限维。若 \(u\in H\)、\(u_V\in V\) 分别解全空间与限制测试方程，证明
\[\lVert u-u_V\rVert_H\le\frac M\alpha\inf_{v\in V}\lVert u-v\rVert_H.\]
说明非对称情形为何不能直接引用正交投影的勾股恒等式。
\end{exercise}
\begin{solution}
限制型仍强制，故 \(u_V\) 存在唯一。误差 \(e=u-u_V\) 满足 \(a(e,w)=0\) 对所有 \(w\in V\)。于是对 \(v\in V\)，
\[\alpha\lVert e\rVert^2\le a(e,e)=a(e,u-v)+a(e,v-u_V)=a(e,u-v)
 \le M\lVert e\rVert\lVert u-v\rVert.\]
误差为零时显然，否则约去并取下确界即可。非对称型一般不是内积，测试正交并不等于 Hilbert 内积正交，故正文 Poisson 问题的勾股公式不能直接套用。
\end{solution}

% 自拟：系数在L无穷范数下的稳定性。
\begin{exercise}\label{b2:ex:15-coefficient-stability}
设 \(A_n,A\) 有共同椭圆下界 \(\lambda>0\)，且 \(\lVert A_n-A\rVert_\infty\to0\)。同一 \(f\in H^{-1}\) 对应的零边界解为 \(u_n,u\)。证明 \(u_n\to u\) 于 \(H_0^1\)，并给出定量估计。
\end{exercise}
\begin{solution}
以 \(A_n\) 为左端，两方程相减得 \(\int A_n\nabla(u_n-u)\cdot\nabla v=\int(A-A_n)\nabla u\cdot\nabla v\)。取 \(v=u_n-u\)，利用椭圆性与 Hölder 不等式，得到
\[\lVert\nabla(u_n-u)\rVert_2\le\lambda^{-1}\lVert A_n-A\rVert_\infty\lVert\nabla u\rVert_2
 \le\lambda^{-2}\lVert A_n-A\rVert_\infty\lVert f\rVert_{H^{-1}}.\]
右端趋零，Poincaré 不等式再给出完整 \(H^1\) 收敛。
\end{solution}

% 自拟：Neumann归一化与相容性。
\begin{exercise}\label{b2:ex:15-neumann}
设 \(\Omega\) 为连通有界 Lipschitz 区域，\(f\in L^2(\Omega)\)。使用零均值 Poincaré 不等式，证明 \(\int\nabla u\cdot\nabla v=\int fv\) 对所有 \(v\in H^1(\Omega)\) 可解，当且仅当 \(\int f=0\)。证明指定 \(\int u=0\) 后解唯一。
\end{exercise}
\begin{solution}
取 \(v=1\) 得必要条件。反之在闭空间 \(H=\{v\in H^1:\int v=0\}\) 上，梯度范数因零均值 Poincaré 不等式而等价于 \(H^1\) 范数，故 Lax--Milgram 给出唯一零均值解。对一般测试函数减去平均值即可落入 \(H\)，而右端常数部分因 \(\int f=0\) 消失，故方程对全部 \(H^1\) 成立。任意两个解之差梯度为零，连通性使其为常数，零均值再使该常数为零。
\end{solution}

% 自拟：漂移项的反对称结构。
\begin{exercise}\label{b2:ex:15-divergence-free-drift}
设 \(b\in C^1(\overline\Omega;\mathbb R^d)\)、\(\operatorname{div}b=0\)，\(A\) 有界且一致椭圆。证明零边界问题 \(-\operatorname{div}(A\nabla u)+b\cdot\nabla u=f\) 对任意 \(f\in H^{-1}\) 唯一可解，其能量估计不含 \(\lVert b\rVert_\infty\)。
\end{exercise}
\begin{solution}
对紧支集光滑 \(u\)，有 \(\int(b\cdot\nabla u)u=\tfrac12\int b\cdot\nabla(u^2)=0\)。双线性连续性和 \(H_0^1\) 密度把该等式推广到全部 \(H_0^1\)。所以型虽可能不对称，却满足 \(a(u,u)\ge\lambda\lVert\nabla u\rVert_2^2\)。有界性仍含 \(\lVert b\rVert_\infty\)，但强制常数不含它。Lax--Milgram 给出 \(\lVert\nabla u\rVert_2\le\lambda^{-1}\lVert f\rVert_{H^{-1}}\) 及唯一性。
\end{solution}

% 自拟：不同共振频率的相容条件与全体解。
\begin{optionalexercise}\label{b2:ex:15-resonance}
在 \((0,\pi)\) 上求解零端点问题 \(-u''-4u=\sin x+q\sin2x\)。确定哪些 \(q\in\mathbb R\) 允许弱解，并在可解时写出全部解。
\end{optionalexercise}
\begin{solution}
型对称，齐次核是 \(\operatorname{span}\{\sin2x\}\)。与核配对得相容条件 \(q\int_0^\pi\sin^22x=0\)，即 \(q=0\)。此时特解为 \(-\tfrac13\sin x\)，故全部解为 \(-\tfrac13\sin x+c\sin2x\)。核的刻画可由一维分布方程先推出 \(u''=-4u\in L^2\)，再由常系数常微分方程及端点条件得到，亦可直接使用正弦展开。因此没有遗漏较粗糙的能量解。
\end{solution}
